This appendix gives the parameter-drift analysis summarized in \Cref{sec:analysis-where}: it measures how far each component moves from the base model and shows that the language model does the work, reproducibly.

\textbf{The language model absorbs the signal.} The language model's relative update is essentially \emph{invariant} to whether the projector is also trainable ($7.72\times10^{-3}$ for LM-only versus $7.67\times10^{-3}$ for LM+projector, agreeing to three significant figures), while under joint training the projector moves only $22\%$ of its solo magnitude ($1.36\times10^{-2}$ vs.\ $6.07\times10^{-2}$). The loss splits its gradient across components, but the language-model solution dominates and the two do not compose: LM+projector is the LM solution, not a superposition. The projector-only model's large projector drift, in turn, feeds the frozen LM visual features off its training distribution, the likely source of its POPE/grounding degradation (\Cref{sec:analysis-pheno}). The vision tower is confirmed frozen, its drift at conversion-noise level across all runs.

\textbf{The update is reproducible and head-concentrated.} Within the language model the update is distributed across all decoder layers but modestly concentrated (layer-level Gini $\approx 0.10$--$0.13$ for attention projections and $\approx 0.06$ for MLP) and sharper at the head level (Gini up to $0.29$), including a standout head (layer~11 query projection, head~22) that dominates its layer. Critically, this structure is \emph{reproducible}: between the LM-only and LM+projector runs the top-drifting layers overlap with Jaccard $=1.00$, and the top heads with Jaccard $=1.00$ in $8$ of $10$ analyzed cells. The optimizer reaches the same low-dimensional attractor regardless of the projector, consistent with the localization-head phenomenon~\citep{kang2025your, kim2025interpreting}, here arising from \emph{explicit} spatial supervision rather than unsupervised training. This reproducible, low-rank signature is what makes the objective a natural fit for the LoRA-based instruction-tuning regime of \Cref{sec:regimes}. We caution that drift is a proxy for functional change, and that we have verified reproducibility across configurations but not yet across random seeds.
